\section{Related Work}

\textbf{Lag-1 and smoother baselines.} Paper~7~\cite{paper7online}
introduced the lag-1 baseline that the present detector gates over;
Paper~8~\cite{paper8multi} ruled out multi-window smoothers as a
fix. Together they motivate the adaptive-procedure class that this
paper enters.

\textbf{Self-trajectory evidence.} Paper~9~\cite{paper9selftraj}
showed that Paper~6's high-capacity collapse was strongly
selection-policy-dependent. The present paper inherits the
fixed-trajectory convention from Papers~7 and~8 and does not address
the trajectory-coupling question; doing so would require
re-running adaptive switching under each method's own trajectory,
which is left as future work.

\textbf{Change-point detection.}
The CUSUM detector~\cite{page1954cusum} is the classical alternative
to the magnitude test used here; it accumulates evidence across
windows rather than reacting to a single inter-window shift.
Bayesian online change-point methods are a richer class. Both are
out of scope; the present paper establishes the floor performance
of the simplest deployable detector.

\textbf{Capacity-conditional gating.} The static
$(K \leq 100 \to \textsf{lag1}, K \geq 200 \to \textsf{fixed})$
gate is, to our knowledge, not previously formalised in the
vulnerability-prioritization literature. It emerges naturally from
the Paper~7 results and is included here as the simplest non-adaptive
baseline against which to compare the change-point procedure.

\textbf{Multi-paper sequence.} HygieneBench (Paper~3)~\cite{paper3hygienebench},
HygienePrio (Paper~4)~\cite{paper4hygieneprio}, Temporal Stability
(Paper~5)~\cite{paper5temporal}, Capacity-Indexed Decay
(Paper~6)~\cite{paper6capacity}, Online Calibration (Paper~7),
Multi-History Calibration (Paper~8), and Self-Trajectory
Evaluation (Paper~9) define the methodological substrate within
which this paper sits.

\textbf{Policy mandates.} CISA BOD~22-01~\cite{cisa2021bod2201},
23-01~\cite{cisa2022bod2301}, NIST SP~800-40
Rev.~4~\cite{nist2022sp80040r4}, and the 2026 DBIR~\cite{verizon2026dbir}
define the operational scheduling regime; the magnitude-test detector
is auditable in a way that gradient-based or Bayesian alternatives
are not, which matters for federal-deployment compliance.
